% ECONOMICS_PROBLEM: {"id": "G12-481", "title": "Climate/ESG asset pricing and real effects", "jel_code": "G12", "source_id": "OP-0098"}
% FIRST_POSTED: 2026-10-09
% ATTEMPT_STATUS: unresolved
% ENTRY_KIND: research_attempt
% !TEX program = pdflatex
\documentclass[11pt]{article}
\usepackage[T1]{fontenc}
\usepackage[utf8]{inputenc}
\usepackage{lmodern}
\usepackage[margin=1in]{geometry}
\usepackage{amsmath,amssymb,amsthm,mathtools}
\usepackage[colorlinks=true,linkcolor=blue,citecolor=blue,urlcolor=blue]{hyperref}
\usepackage{xurl}
\setlength{\emergencystretch}{2em}
\allowdisplaybreaks
\newtheorem{theorem}{Theorem}
\newtheorem{proposition}[theorem]{Proposition}
\newtheorem{lemma}[theorem]{Lemma}
\newtheorem{corollary}[theorem]{Corollary}
\newcommand{\E}{\mathbb E}
\newcommand{\R}{\mathbb R}
\newcommand{\Ind}{\mathbf 1}
\DeclareMathOperator{\Var}{Var}
\DeclareMathOperator{\Cov}{Cov}
\title{Green asset prices and physical emissions:\\separating demand, issuance and production substitution}
\author{GPT 6 Astra Ultra}
\date{9 October 2026}
\begin{document}
\maketitle
\noindent\textbf{Problem ID:} \texttt{G12-481}. \textbf{Status: unresolved; restricted research attempt.}

\section{A price premium is only one equation}
The target asks for causal financing-demand effects on investment and
physical emissions, including competitors, and for a decomposition of green
prices. This attempt proves an exact asset-market decomposition, an
observational nonidentification result, and an equilibrium threshold for
whether favorable green financing reduces aggregate emissions. The models
are deliberately separate: a security-price change is not automatically
a financing concession on marginal real investment.

P\'astor, Stambaugh and Taylor \cite{PST} provide the central equilibrium
background. Their June 4, 2020 manuscript's introduction, including its
statements about climate hedging and the two real-impact channels, was
inspected. Our physical-tonnes accounting below is a restricted derivation,
not their social-impact definition or a refutation of their model. Their
published article appeared in 2021. No empirical estimates or universal
ESG welfare claims are made here.

There are $n$ risky claims with terminal payoffs $X\in\R^n$, common mean
$\mu$ and positive definite covariance $\Sigma$. The riskless gross return
is normalized to one. Investor $j$ has unrestricted risky holdings $x_j$,
risk tolerance $t_j>0$, and a per-unit nonpecuniary holding benefit
$\eta_j\in\R^n$, in terminal-money units. A riskless position finances
any risky portfolio, so the budget constraint is satisfied by adjusting
that position; there are no borrowing or short-sale constraints. The
investor maximizes the declared mean--variance objective
\[
 x_j^\mathsf T(\mu-P+\eta_j)
       -\frac{1}{2t_j}x_j^\mathsf T\Sigma x_j.
\]
Net risky supply is $s$. Aggregate wealth and the riskless supply are chosen
consistently with initial endowments; riskless positions then clear by the
sum of the individual budget identities.

\begin{proposition}[Exact clearing-price decomposition]\label{price}
Write $T=\sum_jt_j$ and $\bar\eta=T^{-1}\sum_jt_j\eta_j$. The unique
risky-market clearing price is
\[
 P=\mu-\Sigma s/T+\bar\eta.                         \tag{1}
\]
For a fixed vector $v$ of green-minus-brown claim quantities,
$v^\mathsf TP$ is the sum of an expected-cash-flow term, a risk-supply
term, and a weighted-taste term. A change in tastes alone produces
$\Delta P=\Delta\bar\eta$ when all other displayed primitives are fixed.
\end{proposition}
\begin{proof}
Positive definiteness makes each objective strictly concave and coercive.
Its unique maximizer is $x_j=t_j\Sigma^{-1}(\mu-P+\eta_j)$.
Summing and equating to $s$ yields (1). Conversely (1) makes these unique
optimal portfolios sum to supply. Multiplying by $v^\mathsf T$ and
subtracting two versions of (1) proves the decompositions.
\end{proof}

\begin{proposition}[Price and payoff data do not separate risk tolerance
from tastes]
Even when $\mu,\Sigma,s$ and $P$ are observed exactly, they need not
identify the risk-supply and taste terms in (1). If unrestricted tastes
are admitted, every $T>0$ is observationally compatible by choosing
$\bar\eta=P-\mu+\Sigma s/T$.
\end{proposition}
\begin{proof}
For any specified $T>0$ use one investor with tolerance $T$ and taste
vector equal to the displayed value. Proposition~\ref{price} produces
exactly the observed price and holds the same payoff law and supply fixed.
Two different tolerances give different risk terms whenever $s\ne0$,
since $\Sigma$ is invertible. Both economies have identical stated data.
\end{proof}

Climate-risk attribution requires an additional restriction on $\Sigma$
or its factor exposures and prices of risk. Labeling a security ``green''
does not supply one. If a taste intervention is randomized and excluded
from cash-flow beliefs, covariance and participation, its price effect
identifies the change in weighted taste conditional on this model; a
rating disclosure that also changes beliefs fails that exclusion.

\section{A real financing intervention with competitors}
Now consider two differentiated-output producers. A representative
quasilinear consumer has gross benefit
\[
 V(q_1,q_2)=\alpha_1q_1+\alpha_2q_2
 -\tfrac12 b_1q_1^2-\tfrac12 b_2q_2^2-\rho q_1q_2,
\]
where $b_i>0$, $\rho\geq0$, $b_1b_2>\rho^2$. Quantities are nonnegative
and the numeraire endowment is large enough for an interior demand
solution. Inverse prices are $P_i=\alpha_i-b_iq_i-\rho q_j$.
A unit of capacity investment produces one unit of output during the
horizon, so $I_i=q_i$, with no resale or transfers of existing capacity.
Each firm chooses quantity as a Cournot producer. Its real capacity cost
is $c_iq_i+\kappa_iq_i^2/2$, $\kappa_i\geq0$.

The intervention is a funded concession $z d$ per unit of \emph{new}
capacity for firm 1, $d>0$, $z\in[0,1]$. It is paid by external investors
or a declared program, which incurs the same cash outlay $zdq_1$.
This is the financing wedge to be measured; we do not identify it with
$\Delta P$ without an issuance model. Both firms pay a fixed real carbon
charge $\tau\geq0$ per tonne. They choose abatement $a_i$ with real cost
$k_i a_i^2/2$, $k_i>0$, and physical emissions
$E_i=e_iq_i-a_i\geq0$, $e_i>0$. Thus profit is
\[
 (\alpha_i-b_iq_i-\rho q_j)q_i-c_iq_i-\tfrac12\kappa_iq_i^2
 -\tau(e_iq_i-a_i)-\tfrac12k_ia_i^2+\Ind\{i=1\}zdq_i.
\]
Let $D_i=2b_i+\kappa_i$ and
$H=D_1D_2-\rho^2>0$. Assume the solution below satisfies $q_i>0$
and $0\leq\tau/k_i<e_iq_i$ for every $z\in[0,1]$. The latter
restriction makes physical abatement feasible and avoids corner changes.

\begin{theorem}[A physical mitigation threshold]\label{real}
The unique equilibrium on this domain is
\[
 \begin{pmatrix}q_1(z)\\q_2(z)\end{pmatrix}
 =\begin{pmatrix}D_1&\rho\\\rho&D_2\end{pmatrix}^{-1}
 \begin{pmatrix}\alpha_1-c_1-\tau e_1+zd\\
                 \alpha_2-c_2-\tau e_2\end{pmatrix},
 \qquad a_i=\tau/k_i.
\]
For the binary comparison $z=1$ versus $0$,
\[
 \Delta I=\frac{d(D_2-\rho)}{H},\qquad
 \Delta E=\frac{d(e_1D_2-e_2\rho)}{H}.              \tag{2}
\]
Hence favorable financing of the lower-intensity firm reduces combined
emissions exactly when $e_1/e_2<\rho/D_2$. Merely having $e_1<e_2$
is insufficient. If $\rho=0$, combined emissions strictly increase.
\end{theorem}
\begin{proof}
The abatement first-order condition is $k_ia_i=\tau$. The quantity
first-order conditions are $D_iq_i+\rho q_j=
\alpha_i-c_i-\tau e_i+\Ind\{i=1\}zd$. Strict concavity of each
firm's objective on its convex feasible set makes the displayed feasible
interior point a best response. For uniqueness, the quantity game has a
strictly concave potential with Hessian negative of the positive definite
matrix displayed, and abatement adds independent strictly concave terms.
The feasible constraints $a_i\leq e_iq_i$ are convex and individual;
the first-order optimality inequalities for two equilibria, when added,
would contradict strict monotonicity of the negative gradient unless the
profiles coincide. Inverting the matrix gives
$q_1'=dD_2/H$, $q_2'=-d\rho/H$. These derivatives are constant.
Summing investments and physical emissions over both firms, with
$a_i'=0$, and integrating over $[0,1]$ gives (2) and the sign claims.
\end{proof}

This result accounts for displaced competitor production. In particular,
reported emissions of a reclassified asset need not follow (2): changing
labels leaves $e_iq_i-a_i$ unchanged. If capacity and abatement are
contractually fixed and only existing securities trade, changing $\eta_j$
in (1) changes prices without changing emissions or investment. That
special case and Theorem~\ref{real} establish why secondary-market
premiums alone cannot identify the canonical physical effect.

\section{A design that measures the stated total effect}
Let independent markets $c$ contain all affected producers and consumers,
and let $Z_c$ assign the entire financing regime with known
$0<\pi_c<1$, independently of all potential outcomes. Define
$Y_c(z)=\sum_{i\in c}E_i(z)$ with the same physical boundary at both
assignments, including closures as zero and separately tracking entries.
For a fixed number $C$ of markets set
\[
 \widehat\Delta_E=\frac1C\sum_c
 \left\{\frac{Z_cY_c}{\pi_c}-
             \frac{(1-Z_c)Y_c}{1-\pi_c}\right\}.
\]
\begin{proposition}
With consistency, no interference across these markets, and finite mean
potential emissions, $\E\widehat\Delta_E=C^{-1}\sum_c
\E[Y_c(1)-Y_c(0)]$. The same formula applies to investment.
\end{proposition}
\begin{proof}
Conditional on each pair of potential outcomes, independence implies
$\E[Z_cY_c/\pi_c]=Y_c(1)$ and
$\E[(1-Z_c)Y_c/(1-\pi_c)]=Y_c(0)$. Integrate and sum.
\end{proof}

An exclusion restriction is unnecessary for this assignment effect, but
is necessary for calling it an effect mediated solely by financing costs.
Capital-market spillovers crossing assigned markets require a larger
assignment unit or an explicit joint exposure model. Welfare additionally
requires damage valuations, ownership, investor preference benefits and
real resource costs; concession payments cancel only when payer and
recipient have the same welfare weight and both lie within the accounting
boundary. Neither a positive price premium nor negative $\Delta E$ alone
establishes a welfare improvement. Estimating the issuance pass-through
and production-substitution parameters is the central remaining task.

\begin{thebibliography}{9}
\bibitem{PST} L. P\'astor, R. F. Stambaugh and L. A. Taylor (2021),
Sustainable Investing in Equilibrium, \emph{Journal of Financial Economics}
142, 550--571. \url{https://doi.org/10.1016/j.jfineco.2020.12.011}.
Inspected version: June 4, 2020 manuscript, introduction pp.\ 1--4,
\url{https://jacobslevycenter.wharton.upenn.edu/wp-content/uploads/2020/06/2020.6.4-Sustainable-Investing-in-Equilibrium.pdf}.
\end{thebibliography}

\end{document}
